\subsection{The ordinary Eisenstein theorem of Hsieh and the actual residual triple}
\label{sec:kt-hsieh-unitary-scope}

This is a direct applicability audit of a published-method input.
It supplies no new vanishing of a Selmer group.  We use the author
version of M.-L.~Hsieh, \emph{Eisenstein congruence on unitary groups
and Iwasawa main conjectures for CM fields}, dated
16 September 2025, retained at
\url{https://www.math.ntu.edu.tw/~mlhsieh/research/IMCCM.pdf}.

\begin{proposition}[KT-HSIEH-DIRECT-SCOPE]
Let $p>3$, let $E/\mathbf Q$ be imaginary quadratic with
$p\nmid\operatorname{disc}(E)$, and let $e$ be even.
The residual representation with semisimplification
\[
       \omega^e\oplus1\oplus\omega^{-e}\quad\hbox{on }G_E
       \tag{HU.1}
\]
does not match the residual Eisenstein branch used in Hsieh's
Theorems 7.18 and 8.14, even after a common scalar-character
twist and a permutation of the three constituents.
\end{proposition}
\begin{proof}
The source assumes that $E$ is $p$-ordinary: for $E/\mathbf Q$
this requires $p$ to split.  Hence there is no direct application
when $p$ is inert.

Suppose $p$ splits and write its chosen ordinary place as $v$,
with conjugate place $\bar v$.  Denote the source's finite branch
character by $\xi$, to distinguish it from our $\psi=\omega^e$.
Hypothesis 7.18(2), printed p.~91, requires that $\xi$ be
unramified at $\bar v$; it also imposes a condition at $v$,
which is not needed for this obstruction.  The exact normalized
three-character parameter in its equation (8.2), printed p.~95,
has residual constituents
\[
       \bar\xi^{-c}\omega,\qquad 1,\qquad
       \bar\xi\omega^{-1}.
       \tag{HU.2}
\]
On $I_v$ the ratio of the first two is therefore $\omega$,
because $\xi^c$ is unramified there.

By contrast, every ratio of two constituents in (HU.1)
restricts to a power of $\omega$ with exponent
$0,\pm e,\pm2e$.  All these exponents are even.
Since $\omega(I_v)=\mathbf F_p^\times$ has even order $p-1$,
none is the character $\omega$.  Common scalar twists cancel
in every ratio, and permutations leave the collection of ratios
unchanged.  Thus neither operation removes the mismatch.
The same argument works at whichever split place is selected
for the CM-type.
\end{proof}

There is a second, independent direction-of-inequality issue.
In the same source, Theorem 7.18 constructs a surjection to
$\Lambda/(L^S(\Psi))$ when $L^S(\Psi)$ is a nonunit.
Definition 7.19 and Corollary 7.20, printed pp.~91--92, give
\[
 \operatorname{Eis}(\Psi,S)
 \subset
 \bigl(L_p^S(-1,\Psi_+\tau_{E/\mathbf Q})
                         L_p^S(\Psi,\Sigma)\bigr).
 \tag{HU.3}
\]
The congruence component in that definition is constructed
under the nonunit hypothesis.  This is a production of
congruences from an $L$-value divisor, and does not state
that a unit $L$-value excludes every cuspidal congruence.
Even if one separately defined an appropriate congruence
ideal in the unit case, the displayed containment would
only put that ideal inside $\Lambda$.
The subsequent Selmer inequalities likewise have the
direction ``Selmer length at least congruence valuation
at least $L$-value valuation.''  A zero lower bound
does not force Selmer vanishing.

The full proof of Theorem 7.18, Definition 7.19,
Corollary 7.20, the exact parameter and GMA setup on
pp.~95--96, and the full proof of Theorem 8.14 on
pp.~101--102 were read for this audit.  The underlying
Hida theory, Fourier--Jacobi nonvanishing and Galois
attachment results are imported source inputs, not
newly reconstructed proofs.  The obstruction above
concerns this precise branch and these theorem hypotheses;
it is not a claim that every possible unitary Eisenstein
comparison has the same restriction.

Consequently this theorem does not provide the missing
integral comparison for our odd off-diagonal module,
nor does it supply a necessary three-unit condition for
the surviving primitive rank-three source.
